% CHECKPOINT35: complete formalization of the frozen checkpoint34 first-step proof.
% Mathematical predecessor SHA256 9ccc6406ec9558d4715ed7837feffb844af40606b2743d6574eb6ef399eba0b5.
\statement{A.114}{Proved}{The next modular central refinement and a full peripheral lift}
Assume the original negative involution belongs to full \(\widehat{\mathrm{GT}}\),
and retain its canonical \(u\), actual \(d=1+2k\), and ordered \(H_d\).
Put \(N_1=\overline{[\ker q,F]}\) and \(M_B=\overline{[N_1,B]}\).
Then \(N_1/M_B\) is pro-two, the exact coordinate criterion
\[
 M_B=N_1\cap\pi_2^{-1}\bigl(\overline{[M_{21},\overline B]}\bigr)
\]
holds, and \(H_d\in M_B\). The original assignments
\(s\mapsto u^{-1}s\), \(\delta\mapsto\delta^d\) define an
actual automorphism of \(B/M_B\), conjugating the original involution
to complex conjugation there. They also admit a full peripheral
lift inducing the specified automorphism of \(B/N_1\), with literal
cusp image \(\delta^d\). This does not identify that lift with the
canonical substitution on the full free group.
\begin{proof}
\medskip\noindent\emph{1. Objects and conclusions}.

Retain \(F=\widehat F(x,y)\), \(B=C_2(s)\amalg C_3(r)\), the actual pure kernel \(F=\ker(B\to S_3)\), and the full signed map \(q:F\to C\). Put
\[
N=\ker q,\quad N_1=\overline{[N,F]},\quad
E_1=F/N_1,\quad B_1=B/N_1.
\]
All closed commutators and quotients below are actual ambient ones. Put
\[
P=F^{(2)},\quad O=\ker(F\to P),\quad
M_2=\ker(P\to C_{(2)}),\quad M_{21}=\overline{[M_2,P]},
\]
\[
\overline B=B/O,\qquad B_2^*=\overline B/M_{21},\qquad
M_B=\overline{[N_1,B]}.
\]
The use of \(\overline B\), whose pure subgroup is the full free pro-two P, is essential. It is not B with a different kernel silently substituted.

We prove:
\[
H^2_{\rm cont}(B_1,\mathbb Q_\ell/\mathbb Z_\ell)=0
\quad(\ell\text{ odd}),                                      \tag{A.114.1}
\]
\[
N_1/M_B\text{ is pro-two},\qquad
M_B=N_1\cap\pi_2^{-1}\bigl(\overline{[M_{21},\overline B]}\bigr),
\qquad H_d\in M_B.                                            \tag{A.114.2}
\]
Here \(\pi_2\) on \(N_1\) is the restriction of the actual map B\(\to\)\(\overline B\); it agrees with the maximal pro-two map on F.

Consequently the original assignments
\[
s\mapsto v=u^{-1}s,\qquad \delta=sr\mapsto\delta^d
\]
give a genuine automorphism of B/\(M_B\), with positive or negative r-factor frame as in A.108 and with those frames lifting the actual ones on \(B_1\). The closed-normal central lifting criterion of A.3/A.34 then yields a full peripheral automorphism of B inducing the specified \(\Psi_1\) on \(B_1\). Its cusp can be normalized literally to \(\delta^d\).

This is a lift of the quotient automorphism. It is not an identification of that lift's full pure substitution with the actual canonical \(\psi_d\), nor a full conjugacy between the original \(\alpha\) and \(c_B\).

\medskip\noindent\emph{2. The actual pure product and the two-primary comparison}.

Write \(C_o=\prod_{p\ {\rm odd}}\operatorname{SL}_2(\mathbb Z_p)\).
Beyl's Proposition5.7(a) gives \((C_o)_{\rm ab}=C_3\), supported at three, with the original character \(\chi(X)=\chi(Y)=1\). Thus \(C_o\) has no nontrivial pro-two quotient: any nontrivial finite two-group quotient has a nontrivial \(C_2\) character, which would contradict this abelianization.

The maximal pro-two quotient of the actual \(C=C_{(2)}\times C_o\) is therefore \(C_{(2)}\). Right exactness of the maximal pro-two quotient gives
\[
\pi_2(N)=M_2,\qquad \pi_2(N_1)=M_{21}.                         \tag{A.114.3}
\]
For clarity, \(P/\pi_2(N)\) is pro-two and has exactly the universal property of the maximal pro-two quotient of F/N. This proves the first equality; continuous commutator functoriality and compactness give the second.

Define
\[
P^*=P/M_{21},\qquad
D_o=(F_{\rm ab})_o\times_{(C_o)_{\rm ab}}C_o.                   \tag{A.114.4}
\]
The map from F to \(D_o\) is onto. First match the odd abelian vector. Its discrepancy from the requested \(C_o\)-coordinate lies in \(C_o'\), and is corrected by an F' lift, because \(q_o(F')=C_o'\). The map to \(P^*\) is onto and kills \(N_1\) by (A.114.3). A.106's exact coordinate equality proves that the combined map induces an injection
\[
E_1\hookrightarrow P^*\times D_o.                             \tag{A.114.5}
\]
Indeed, a zero \(P^*\)-coordinate gives zero signed dyadic coordinate and zero dyadic abelian vector; the zero \(D_o\)-coordinate gives zero odd signed and abelian coordinates. The A.106 criterion then gives membership in \(N_1\).

The kernel of \(D_o\to C_o\) is pro-odd abelian. Hence any map from \(D_o\) to a finite two-group kills that kernel and factors through \(C_o\), and is trivial. Profinite Goursat applied to the subdirect image in (A.114.5) now yields the exact product
\[
E_1=P^*\times D_o.                                            \tag{A.114.6}
\]
Any common quotient would be pro-two because \(P^*\) is; the preceding argument excludes a nontrivial such quotient of \(D_o\).

The factor \(D_o\) is normal in \(B_1\). More intrinsically it is the kernel of
\[
B_1\longrightarrow B/(ON_1)=B_2^*.
\]
Both O and \(N_1\) are normal in B: O is characteristic in the actual F, and A.107 proves the required full signed geometric stability of \(N_1\). The image of \(N_1\) under B\(\to\)\(\overline B\) is exactly M21 by (A.114.3). Thus this quotient identification, including its \(S_3\) action, is literal.

Let \(\mathcal A_2=\mathbb Q_2/\mathbb Z_2\), always with trivial action. The pro-odd central kernel in \(D_o\to C_o\) has no positive continuous cohomology with two-primary coefficients, by finite-quotient averaging. Beyl Proposition~5.7(b), the complete product argument of A.106, and the odd abelianizations give
\[
H^1(D_o,\mathcal A_2)=H^2(D_o,\mathcal A_2)=0.                 \tag{A.114.7}
\]
The low-degree Lyndon--Hochschild--Serre spectral sequence for \(D_o\to B_1\to B_2^*\) therefore gives
\[
H^2(B_2^*,\mathcal A_2)\ \xrightarrow{\ \mathrm{inf}\ }\ 
H^2(B_1,\mathcal A_2)
\quad\text{is an isomorphism}.                               \tag{A.114.8}
\]
No assertion that either side is zero is used.

\medskip\noindent\emph{3. The odd cohomology calculation, including the prime-three lattice}.

Set A=N/\(N_1\). It is central in E1. Its map to \(F_{\rm ab}\) has kernel
\[
K=(N\cap F')/N_1,
\]
which is pro-two by A.106, and image
\[
L=\ker(F_{\rm ab}\to C_{\rm ab}).                              \tag{A.114.9}
\]
The image assertion follows from right exactness of abelianization: the image of compact N is closed. Since A is profinite abelian, its odd primary components are canonically those of L. Consequently
\[
A_\ell\simeq\mathbb Z_\ell^2\quad(\ell\ge5),\qquad
A_3\simeq L_3=\{(a,b)\in\mathbb Z_3^2:a+b\equiv0\bmod3\}.     \tag{A.114.10}
\]
Their actual \(S_3\) actions are the restrictions of the original swap \(\iota\) and cyclic matrix \(\rho\) on \(F_{\rm ab}\).

Fix an odd prime \(\ell\), and put \(\mathcal A_\ell=\mathbb Q_\ell/\mathbb Z_\ell\). A.106's actual product/Beyl argument gives
\[
H^2(C,\mathcal A_\ell)=0.
\]
Apply LHS to the central extension \(A\to E_1\to C\). The action of C on
\[
M_\ell=H^1(A,\mathcal A_\ell)
\]
is trivial. In total degree2 its possible associated graded terms are:
\begin{itemize}
\item the quotient of \(H^2(C,\mathcal A_\ell)=0\);
\item a \emph{submodule} of \(H^1(C,M_\ell)\);
\item a \emph{submodule} of \(H^2(A,\mathcal A_\ell)\).

\end{itemize}
The submodule qualification is essential: for the positions (1,1) and (0,2) there are no incoming differentials in this first-quadrant spectral sequence. Thus it is legitimate to use the vanishing of their \(S_3\) invariants. The canonical filtration and differentials are equivariant under conjugation by \(B_1\); inner E1 actions act trivially on these cohomology groups, so the action factors through the actual \(S_3\).

For \(\ell\ne3\), \(H^1(C,M_\ell)=0\). For \(\ell=3\), it is
\[
\operatorname{Hom}(C_{\rm ab},M_3)=M_3[3].
\]
The actual \(S_3\) action on the \(C_3\) abelian quotient of C is trivial: the swap preserves the character sum, and the cyclic substitution sends each of x,y,z to a generator with that same character modulo3. The central cyclic sign has character zero. Hence the invariants here are the invariants of the dual root lattice. They vanish. Explicitly, with
\[
a=(1,-1),\qquad b=(0,3)
\]
as a basis of L3,
\[
(\iota-1)a=-2a,\qquad(\rho-1)a=b.
\]
Since2 is a unit in \(\mathbb Z_3\), these relations kill both basis vectors in the coinvariants. Thus \(L_{3,S_3}=0\) and \(M_3[3]^{S_3}=0\).

Finally \(H^2(A,\mathcal A_\ell)\) is the continuous alternating-form module on \(\Lambda^2 A_\ell\). This follows from the two procyclic factors of the rank-two lattice: each has cohomological dimension one, and the product cross term is the alternating bicharacter. Other primary factors have zero positive \(\ell\)-cohomology. The swap has determinant\(-\)1 on this lattice, so acts by\(-\)1 on its second cohomology. Since \(\ell\) is odd, its invariant subgroup is zero.

Each associated graded term therefore has zero \(S_3\) invariants. The finite filtration implies
\[
H^2(E_1,\mathcal A_\ell)^{S_3}=0.                             \tag{A.114.11}
\]

Now use \(E_1\to B_1\to S_3\). We have
\[
H^1(E_1,\mathcal A_\ell)=
\operatorname{Hom}(F_{\rm ab},\mathcal A_\ell)
=\mathcal A_\ell^2
\]
with the actual standard dual action. The remaining total-degree2 terms vanish:
\[
H^2(S_3,\mathcal A_\ell)=0,\qquad
H^1(S_3,\mathcal A_\ell^2)=0.                                 \tag{A.114.12}
\]
For \(\ell\ge5\) this is finite-group averaging. At three use the normal \(C_3\) and the quotient \(C_2\). On the standard dual, the norm \(1+\rho+\rho^2\) is zero, while \(\rho-1\) is onto the divisible module: its determinant is three, and the adjugate formula plus divisibility supplies preimages. Thus its \(C_3\) first cohomology is zero; \(C_2\) has invertible order. For the trivial coefficient module, \(H^2(C_3,\mathcal A_3)=\mathcal A_3/3\mathcal A_3=0\), and positive \(C_2\) cohomology vanishes. These prove (A.114.12). Equations (A.114.11)--(A.114.12) prove (A.114.1).

\medskip\noindent\emph{4. The actual free-pro-two modular comparison}.

The quotient \(\overline B=B/O\) has pure subgroup P, a full free pro-two group of rank2, and quotient \(S_3\) with the actual standard action on the pure abelianization.

Its continuous two-primary \(H^2\) is zero. Indeed \(H^2(P,\mathcal A_2)=0\). For \(H^1(P,\mathcal A_2)=\mathcal A_2^2\), the \(C_3\) invariants vanish because \(\rho-1\) has determinant3, a two-adic unit; higher \(C_3\) cohomology vanishes by averaging. Therefore
\[
H^1(S_3,\mathcal A_2^2)=0.
\]
For the trivial coefficients, normal-\(C_3\) averaging reduces
\[
H^2(S_3,\mathcal A_2)=H^2(C_2,\mathcal A_2)
=\mathcal A_2/2\mathcal A_2=0.
\]
The same LHS total-degree2 calculation gives
\[
H^2(\overline B,\mathcal A_2)=0.                              \tag{A.114.13}
\]
This direct proof avoids any unverified comparison of extension classes.

There is the actual commutative extension diagram
\[
\begin{array}{ccccccccc}
1&\to&N_1&\to&B&\to&B_1&\to&1\\
 &&\downarrow&&\downarrow&&\downarrow\\
1&\to&M_{21}&\to&\overline B&\to&B_2^*&\to&1 .
\end{array}                                                  \tag{A.114.14}
\]
Use the same transgression convention in both five-term sequences. Equations (A.114.8),(A.114.13) imply that every B-invariant continuous character \(\chi:N_1\to\mathcal A_2\) has the form
\[
\chi=\chi_2\circ\pi_2|_{N_1}+\lambda|_{N_1},                  \tag{A.114.15}
\]
where \(\chi_2:M_{21}\to\mathcal A_2\) is \(\overline B\)-invariant and \(\lambda:B\to\mathcal A_2\) is a character. To obtain this, lift the transgression of \(\chi\) through (A.114.8), choose a character with that lower transgression using (A.114.13), and use naturality and exactness on the difference. No surjectivity of a false full-kernel map is assumed; the required actual map was proved in (A.114.3).

For odd-primary characters, (A.114.1) makes the transgression zero, so every B-invariant \(\chi:N_1\to\mathcal A_\ell\) extends to a character of B. But \(N_1\subset F'\subset B'\), so all such characters are zero. Character separation of the compact abelian group \(N_1/M_B\) proves it is pro-two.

Commutator functoriality gives the forward inclusion in the coordinate formula (A.114.2). Conversely, for an element of its right side, every character in (A.114.15) vanishes: the first term kills \([M_{21},\overline B]\), and the second kills \(N_1\subset B'\). The odd-primary characters already vanish. Invariant-character separation proves the reverse inclusion.

The actual full canonical hexagon belongs to \(N_1\) by A.106, and its whole pro-two image is one by A.89.3 with A.42/A.45 for this same u,d. Therefore the exact criterion gives \(H_d\in M_B\). This is stronger than a dyadic coefficient equality. In fact (A.114.3) and (A.114.2) also give the actual compact isomorphism
\[
N_1/M_B\ \simeq\ 
M_{21}/\overline{[M_{21},\overline B]}.
\]
No claim that \(N_1\) itself is pro-two is made; A.110 proves the contrary.

\medskip\noindent\emph{5. The genuine next modular map and its factor frames}.

The ordered identity
\[
H_d=u((v\delta^d)^3)^{-1}u^{-1}
\]
now kills the cubic in B/\(M_B\). Since v\textsuperscript{2}=1, the free profinite product defines a homomorphism B\(\to\)B/\(M_B\) with s\(\to\)v and r\(\to\)v\(\delta^d\). Its pure restriction is the original \(\psi_d\).

This map kills \(M_B\). Indeed \(\psi_d(N_1)\subset N_1\) by A.107, while \(N_1\)/\(M_B\) is central in B/\(M_B\); the images of all commutators [\(N_1\),B] are therefore1. It descends to an endomorphism \(\Psi_M\) of B/\(M_B\). Its pure map is onto because
\[
N_2=\overline{[N_1,F]}\subset M_B
\]
and A.107 already proves \(\psi_d\) is an automorphism of F/\(N_2\). The finitely generated profinite pure quotient is Hopfian, hence this induced pure map is an automorphism. The \(S_3\) map is the identity, so \(\Psi_M\) is an automorphism of B/\(M_B\). The same generator calculation gives \(\alpha\Psi_M=\Psi_M c_B\).

Lift the actual A.108 frame h to \(\widetilde h\in F/M_B\). The two elements
\[
\Psi_M(r),\qquad \widetilde hR_e\widetilde h^{-1}
\]
have the same image in \(B_1\), so differ by an element of the central pro-two group \(N_1\)/\(M_B\). Both have order dividing three. Their central difference consequently has cube1 and is one. Thus every such lift is the required r-factor frame. The s-factor has its genuine full frame already: v is an involution in B with \(S_3\) image s, so the full finite-factor torsion theorem and the \(S_3\) normalization give \(v=gsg^{-1}\) for a pure g. The cusp is mapped literally to \(\delta^d\). These are actual peripheral-power data at the next central refinement, not a use of a conjugacy-only theorem on powers.

\medskip\noindent\emph{6. Exact multiplier reduction and the full lift it permits}.

Choose an actual Galois element with cyclotomic multiplier d and its original normalized action \(\gamma\) on B. The original LM7\_PUSH E65.8 gives \(\gamma(\delta)=\delta^d\); its original LM7\_PUSH E.65.15 proves, over every \(\mathbb Z_\ell\), signed covariance for the actual \(B_3\) representation. Restriction to the based pure generators \(x=u_{\rm braid}^2\), \(y=v_{\rm braid}^2\) gives
\[
q\gamma=\operatorname{Ad}(\mathcal V_\gamma)q.
\]
Hence \(\gamma(N)=N\), and functoriality gives preservation of \(N_1\) and \(M_B\). This is a use of that explicit original theorem, not an assumption that an arbitrary peripheral automorphism preserves the signed kernel.

Compose \(\Psi_M\) with \(\gamma^{-1}\). The common multiplier cancels, including its r-exponent e. The resulting map on B/\(M_B\) fixes \(\delta\) and preserves each positive peripheral conjugacy class, with specified literal conjugators. It descends to the corresponding automorphism \(\eta_1\) on B/\(N_1\).

The closed-normal modular central criterion in the final paragraph of A.34, proved by the complete A.3 finite-count argument, now applies with relative kernel \emph{\(N_1\)}:
\[
M_B=[N_1,B].
\]
It gives a full positive-peripheral automorphism \(\widetilde\eta\) inducing \(\eta_1\), with the specified conjugators modulo \(N_1\). Its cusp conjugator can be taken in \(N_1\), and the usual inner correction by that element fixes \(\delta\) literally without changing the quotient data. Then
\[
\widetilde\Psi=\gamma\widetilde\eta
\]
is a full peripheral automorphism of B inducing \(\Psi_1\) on \(B_1\) and sending \(\delta\) literally to \(\delta^d\).

The central criterion has been used at exactly the required level. Before the new \(H_d\in[N_1,B]\) argument, it yielded only a lift of the coarser action on B/N. It did not directly lift \(\Psi_1\).

\end{proof}
\noindent\emph{Dependencies and scope.}
The proof uses the actual signed marking of A.103--A.108,
the complete A.3/A.34 closed-normal modular lifting criterion,
and the original \(\mathrm{LM7\_PUSH}\) E.65.8/E.65.15 cusp and covariance
formulas. Its continuous cohomological input is Beyl's Proposition~5.7(a,b),
with the full product and lattice calculation printed above or in A.114.
The original full involution and coherent frame problems remain open.
